\chapter{Referências normativas e bibliografia}
\label{apD}

\begin{objetivos}
Ao final deste apêndice você deve ser capaz de:
\begin{itemize}
  \item citar uma norma com edição e ano, e reconhecer uma citação incompleta;
  \item distinguir obra-fonte, bibliografia de apoio e documentação de fabricante;
  \item identificar qual referência sustenta cada afirmação técnica do texto;
  \item montar a lista de referências do seu próprio projeto.
\end{itemize}
\end{objetivos}

\section{Como as referências são citadas nesta obra}

Três regras valem para todo o texto:

\begin{enumerate}
  \item \textbf{toda afirmação normativa cita a edição e o ano}. ``A norma define'' não é referência:
        ``a IEC 61131-3, 4ª edição, de maio de 2025'' é;
  \item \textbf{nada é reproduzido das normas}. Elas foram consultadas para conferência factual, e o
        texto desta apostila não transcreve trechos, não copia figuras e não inventa numeração de
        seção. Quem precisa do texto normativo o obtém junto ao organismo que o publica;
  \item \textbf{figura ou tabela de terceiro entra com crédito}. Esquema de catálogo, manual de
        fabricante ou figura de outra obra traz a indicação de fonte --- e, sempre que possível, é
        redesenhado em vez de copiado.
\end{enumerate}

\index{referências normativas}\index{bibliografia}

\section{Obras de origem desta apostila}

Esta apostila foi produzida a partir de duas obras do mesmo grupo de autores, incorporadas e
atualizadas:

\begin{itemize}[leftmargin=1.2cm, label={}]
  \item TAVARES, M.Sc. Leonardo de Oliveira; MONTEIRO, M.Sc. Natália Nogueira. \emph{Controladores
        Lógicos Programáveis --- apostila final (teórica e prática)}. 1ª edição, Rev. 0, 2017
        (revisão de 2021, Rev. 05D). Instituto Federal Fluminense, Campus Campos dos Goytacazes.
        \textbf{Obra-fonte principal:} fundamentos, hardware, memória, cartões de E/S, princípio de
        funcionamento, linguagens, família IEC 61131 e a bateria de práticas.
  \item VIANNA, D.Sc. William da Silva. \emph{Controlador Lógico Programável}. Instituto Federal Fluminense,
        Campus Campos dos Goytacazes, 2023. \textbf{Obra-fonte complementar:} características,
        estrutura básica, mapa de memória, critérios de especificação e compra, e noções de sistema
        SCADA com uso de CLP.
\end{itemize}

\section{Normas e documentos normativos citados}

A Tabela~\ref{tab:normas} reúne as normas mencionadas no texto. A coluna de edição e ano registra o
que foi \textbf{efetivamente conferido em fonte pública}; onde a edição não pôde ser confirmada, a
célula indica que ela deve ser conferida no catálogo do organismo, e o texto desta apostila não faz
afirmação técnica que dependa dela.

\begin{longtable}{p{3.6cm} p{5.0cm} p{6.4cm}}
  \caption{Normas citadas, assunto e edição conferida.}
  \label{tab:normas} \\
  \toprule
  Norma & Assunto & Edição e ano \\
  \midrule
  \endfirsthead

  \multicolumn{3}{l}{\small\itshape Continuação da Tabela~\ref{tab:normas}.} \\
  \toprule
  Norma & Assunto & Edição e ano \\
  \midrule
  \endhead

  \midrule
  \multicolumn{3}{r}{\small\itshape Continua na página seguinte.} \\
  \endfoot

  \bottomrule
  \endlastfoot

  IEC 61131-1 & Informações gerais sobre controladores programáveis &
    \emph{Conferir no catálogo}; assunto citado no Capítulo~\ref{cap:familia-iec} \\
  IEC 61131-2 & Requisitos de equipamento e ensaios &
    \emph{Conferir no catálogo}; assunto citado nos Capítulos~\ref{cap:es} e~\ref{cap:familia-iec} \\
  IEC 61131-3 & Linguagens de programação &
    1ª ed.\ 1993; 2ª ed.\ 2003; 3ª ed.\ 2013 (POUs orientadas a objeto; deprecia a Lista de
    Instruções); \textbf{4ª ed.\ maio de 2025 (vigente; removeu a Lista de Instruções)} \\
  IEC 61131-4 & Orientações ao usuário &
    \emph{Conferir no catálogo}; assunto citado no Capítulo~\ref{cap:familia-iec} \\
  IEC 61131-5 & Comunicações &
    \emph{Conferir no catálogo}; assunto citado no Capítulo~\ref{cap:familia-iec} \\
  IEC 61131-6 & Segurança funcional &
    \emph{Conferir no catálogo}; assunto citado no Capítulo~\ref{cap:familia-iec} \\
  IEC 61131-7 & Programação de controle \emph{fuzzy} &
    \emph{Conferir no catálogo}; assunto citado no Capítulo~\ref{cap:familia-iec} \\
  IEC 61131-8 & Implementação e aplicação das linguagens &
    \emph{Conferir no catálogo}; assunto citado no Capítulo~\ref{cap:familia-iec} \\
  IEC 61131-9 & Interface digital de comunicação ponto a ponto (SDCI, comercializada como IO-Link) &
    2013; assunto citado nos Capítulos~\ref{cap:familia-iec} e~\ref{cap:redes-scada} \\
  IEC 61131-10 & Formato de intercâmbio XML de projetos IEC 61131-3, derivado dos esquemas da PLCopen &
    \emph{Conferir no catálogo}; assunto citado nos Capítulos~\ref{cap:familia-iec}
    e~\ref{cap:interoperabilidade} \\
  IEC 61499 (partes 1, 2 e 4) & Blocos de função para medição e controle de processo distribuído &
    Publicada em 2005; parte 1, 2ª ed.\ 2012; parte 2, 2ª ed.\ 2012; parte 4, 2ª ed.\ 2013. A parte 3
    foi retirada em 2008 \\
  IEC 61508 & Segurança funcional de sistemas elétricos, eletrônicos e eletrônicos programáveis
    (E/E/PE) &
    1ª ed.\ 1998--2000; \textbf{2ª ed.\ 2010} \\
  IEC 61511 & Segurança funcional para a indústria de processo & \emph{Conferir no catálogo};
    aplicação setorial da IEC 61508, citada no Capítulo~\ref{cap:seguranca} \\
  IEC 62443-1-1 & Terminologia, conceitos e modelos da série de segurança para automação &
    2009 (publicação da série conjunta ISA/IEC) \\
  IEC 62443-2-1 & Programa de segurança do sistema de automação e controle & 2024 \\
  IEC 62443-2-4 & Requisitos para prestadores de serviço de integração e manutenção & 2023 \\
  IEC 62443-3-2 & Avaliação de risco de segurança para projeto de sistema & 2020 \\
  IEC 62443-3-3 & Requisitos de segurança de sistema e níveis de segurança & 2013 \\
  IEC 62443-4-1 & Requisitos do ciclo de vida de desenvolvimento seguro & 2018 \\
  IEC 62443-4-2 & Requisitos técnicos de segurança para componentes & 2019 \\
  ISO/IEC 27001 & Sistemas de gestão da segurança da informação & 3ª ed.\ outubro de 2022; emenda
    1 de 2024 \\
  ISO/IEC/IEEE 60559 & Aritmética de ponto flutuante & 2020, conteúdo idêntico ao IEEE 754-2019 \\
  IEEE 754 & Aritmética de ponto flutuante binária e decimal & \textbf{2019}, substituindo a edição de
    2008 \\
  NIST Cybersecurity Framework & Referência de programa de segurança cibernética & Versão 2.0, 2024 \\
  ISO/IEC 17065 e ISO/IEC 17025 & Avaliação da conformidade e competência de laboratórios de ensaio &
    \emph{Conferir no catálogo}; citadas no contexto de certificação de produto
    (Capítulo~\ref{cap:seguranca}) \\
\end{longtable}

\paragraph{Sobre as edições.} As séries IEC são revisadas continuamente: uma parte pode ser confirmada,
atualizada ou retirada. Toda edição indicada acima corresponde à publicação consultada e registrada em
2026, e \textbf{não} é garantia de que a versão vigente seja a mesma no futuro. Antes de citar uma
norma em laudo, contrato ou trabalho acadêmico, confirme a edição atual no catálogo do organismo
emissor.

\paragraph{Adoções nacionais.} A ABNT publica adoções de normas IEC com a designação \emph{NBR IEC}.
Isso significa que alguns dos documentos acima podem existir em versão brasileira, com a mesma
numeração de parte. Confirme no catálogo da ABNT qual parte está disponível e em que edição --- e
nunca utilize cópia não autorizada de norma.

\section{Bibliografia consolidada}

As obras abaixo reúnem as bibliografias das duas apostilas de origem, com duplicatas consolidadas e
grafias evidentemente incorretas corrigidas. A marcação indica se a obra permanece uma referência
útil ao estudante atual.

\subsection{Livros e apostilas}

\begin{itemize}[leftmargin=1.2cm, label={}]
  \item FRANCHI, C. M.; CAMARGO, V. L. \emph{Controladores Lógicos Programáveis --- Sistemas
        Discretos}. 2ª ed. São Paulo: Érica, 2010. \textbf{Atual.} Referência de base, ainda alinhada
        à segunda edição da norma; confira a terminologia nas edições posteriores.
  \item PRUDENTE, F. \emph{Automação Industrial --- PLC: Teoria e Aplicações}. 2ª ed. Rio de Janeiro:
        LTC, 2011. \textbf{Atual.} Tratamento didático de Ladder e de comandos, útil para quem parte
        do comando elétrico.
  \item NATALE, F. \emph{Automação Industrial}. São Paulo: Érica, 1992. \textbf{Histórica.} Cobre o
        período anterior à difusão das redes industriais; mantida por registro e para consulta a
        sistemas antigos.
  \item OLIVEIRA, J. C. P. \emph{Controlador Programável}. São Paulo: Makron Books, 1993.
        \textbf{Histórica.}
  \item MORAES, C. M.; CASTRUCCI, P. L. \emph{Engenharia de Automação Industrial}. 2ª ed. Rio de
        Janeiro: LTC, 2007. \textbf{Atual} para instrumentação, arquitetura de automação e sistemas
        industriais.
  \item OIKO, R. K.; SARROUF, L. P. Controladores Programáveis --- como comprar. \emph{Instec}, São
        Paulo, n.~52, p.~32--45, março de 1992. \textbf{Histórica.} Mantida por ser registro do
        problema de especificação, tema retomado no Capítulo~\ref{cap:porte}.
  \item SOUZA, F. M. \emph{Instrumentação --- Automação Básica}. Espírito Santo, 1999.
        \textbf{Histórica.}
  \item VIANNA, W. S. \emph{Apostila de Controlador Lógico Programável}. Centro Federal de Educação
        Tecnológica de Campos, 1995. \textbf{Histórica.} Primeira versão da linha que deu origem à
        obra-fonte complementar.
  \item VIANNA, W. d. S. \emph{Controlador Lógico Programável}. Campos dos Goytacazes, 2008.
        \textbf{Histórica.}
\end{itemize}

\subsection{Artigos, relatórios e documentos técnicos}

\begin{itemize}[leftmargin=1.2cm, label={}]
  \item CHOULIARAS, V. A.; NUNEZ-YANEZ, J. L. An IEEE 754 floating point engine designed with an
        electronic system level methodology. \emph{IEEE}, 2007. \textbf{Pontual.} Sustenta a
        representação em ponto flutuante tratada no Capítulo~\ref{cap:memoria}.
  \item HWANG, Y.-M.; KIM, M. G.; RHO, J.-J. Understanding Internet of Things (IoT) diffusion:
        focusing on value configuration of RFID and sensors in business cases (2008--2012).
        \emph{Information Development}, v.~32, n.~4, 2016. \textbf{Pontual.}
  \item LÜCKENHAUS, M. Machine Vision in IIoT. \emph{Quality Magazine}, 2016. \textbf{Pontual.}
        Pano de fundo para a visão computacional tratada no Capítulo~\ref{cap:ti-to}.
  \item RAY, T.; SURESH, A.; DASH, P. S.; BANERJEE, P. K. Instrumentation and automation system at
        coal leaching pilot plant. \emph{International Journal of Automation and Control}, v.~9,
        p.~158--171, 2015. \textbf{Pontual.}
  \item SASAJIMA, H.; ISHIKUMA, T.; HAYASHI, H. Future IIoT in process automation --- latest trends
        of standardization in industrial automation, IEC/TC65, 2015. \textbf{Pontual.} Origem da
        discussão de IIoT e padronização retomada nos Capítulos~\ref{cap:iiot}
        e~\ref{cap:ti-to}.
  \item ROCKWELL AUTOMATION. \emph{Controladores Logix5000 --- Manual de Referência Geral do Conjunto
        de Instruções}. 2001. \textbf{Fonte do dialeto.} É a origem dos mnemônicos \texttt{XIC},
        \texttt{XIO}, \texttt{OTL} e \texttt{OTU} documentados no Capítulo~\ref{cap:iec-versus-legado}
        e no Apêndice~\ref{apB}. \textbf{Não é norma} e não deve ser usada como referência de
        linguagem padronizada.
  \item EESC/USP. \emph{Memórias}. Grupo de Sistemas Digitais, s.d. \textbf{Histórica.}
  \item FEEC/UNICAMP. \emph{Representação Numérica}. s.d. \textbf{Histórica.}
  \item \emph{The father of invention}: Dick Morley looks back on the 40th anniversary of the PLC.
        \emph{Manufacturing Automation}, s.d. \textbf{Pontual.} Fonte do relato sobre a origem do
        controlador, usado no Capítulo~\ref{cap:evolucao}.
\end{itemize}

\subsection{Documentação técnica e fontes de organismos}

\begin{itemize}[leftmargin=1.2cm, label={}]
  \item OPC FOUNDATION. \emph{OPC Unified Architecture}. Documentação técnica oficial, disponível
        publicamente. Consultada para os Capítulos~\ref{cap:familia-iec} e~\ref{cap:iiot}.
  \item OASIS. \emph{MQTT Version 5.0} (e versão 3.1.1). Especificação padrão, disponível
        publicamente. Consultada para o Capítulo~\ref{cap:iiot}.
  \item ECLIPSE FOUNDATION. \emph{Eclipse Sparkplug Specification}, versão 3.0. Consultada para o
        Capítulo~\ref{cap:iiot}.
  \item MODBUS ORGANIZATION. \emph{Modbus Application Protocol Specification}. Consultada para o
        Capítulo~\ref{cap:redes-scada}.
  \item PLCopen. \emph{Biblioteca padronizada de blocos de função} e esquemas XML de intercâmbio.
        Consultada para os Capítulos~\ref{cap:familia-iec} e~\ref{cap:interoperabilidade}.
  \item NIST. \emph{National Vulnerability Database} e \emph{Cybersecurity Framework 2.0}.
        Consultados para o Capítulo~\ref{cap:seguranca}.
  \item CISA. \emph{Known Exploited Vulnerabilities Catalog} e avisos ICS-CERT. Consultados para o
        Capítulo~\ref{cap:seguranca}.
  \item Documentação técnica dos fabricantes de controladores: manuais de instruções, de comunicação e
        de diagnóstico. São a fonte primária para endereçamento, blocos de sistema e comportamento
        específico de cada ambiente --- e \textbf{não} devem ser citados como norma.
\end{itemize}

\begin{atencao}
Uma referência antiga não é uma referência errada: é uma referência \emph{datada}. O erro está em citar
como atual o que descreve outra geração de equipamento --- por exemplo, apresentar a segunda edição da
IEC 61131-3 como vigente, ou explicar endereçamento por área de memória como se fosse elemento da
norma. Ao usar qualquer obra desta lista, confira em que edição da norma ela se apoia.
\end{atencao}

\begin{pratica}
Monte a lista de referências do projeto que você desenvolveu nas práticas do Apêndice~\ref{apA}:

(a) separe o que é norma, o que é bibliografia de apoio e o que é documentação de fabricante;
(b) para cada norma citada, registre a edição e o ano que você efetivamente consultou;
(c) identifique uma afirmação técnica do seu projeto que \emph{não} tem referência e providencie uma;
(d) substitua qualquer citação de manual de fabricante usada como se fosse norma;
(e) escreva duas linhas explicando por que a data de consulta de uma norma faz diferença em um
projeto que terá de ser mantido por dez anos.
\end{pratica}

\begin{exercicios}
\begin{enumerate}[label=\alph*)]
  \item Explique por que a obra diz que ``a norma define'' não é uma referência aceitável.
  \item Qual é a diferença entre a obra-fonte principal e a obra-fonte complementar desta apostila?
  \item Por que o manual de instruções de um fabricante não pode ser usado como referência de linguagem
        padronizada?
  \item Escolha três normas da Tabela~\ref{tab:normas} e diga em que capítulo cada uma é usada e para
        afirmar o quê.
  \item Explique por que a apostila marca algumas normas como ``conferir no catálogo'' em vez de
        informar uma edição.
  \item Por que a edição vigente da IEC 61131-3 importa para quem vai escrever um programa hoje?
  \item O que significa dizer que a ABNT publica uma norma como ``NBR IEC''?
  \item Por que uma obra marcada como \emph{histórica} nesta bibliografia continua listada?
  \item Cite duas obras desta bibliografia que permanecem úteis ao estudante atual, e justifique.
  \item Por que a apostila afirma que nada foi reproduzido das normas, e por que isso é uma questão de
        direito autoral além de rigor técnico?
\end{enumerate}
\end{exercicios}
